% GENERATED FILE. Edit canonical lessons, then run npm run generate:books.
% canonical-source-hash: fnv1a64:9f26d31d048caa0f
% canonical-lessons: GE-C484-stirn, GE-C484-wange, GE-C484-kinn, GE-C484-augenbraue, GE-C484-wimper

\chapter{Stirn, Wange und Kinn}
\label{ch:ge-stirn-wange-und-kinn}
\hlchaptermodality{484}

\begin{chapteropening}
\textbf{By the end of this chapter:} \emph{I can describe the parts of a face.}

Complete the last lesson of chapter 484: I can describe the parts of a face.
\end{chapteropening}

\section[die Stirn]{die Stirn — forehead}

\label{lesson:GE-C484-stirn}



\begin{quote}
Before the new one: say the German for wheelchair, then the German for to heal.
\end{quote}

\subsection*{You'll want to know: die Stirn}
\textbf{die Stirn} — \textquotedblleft{}forehead\textquotedblright{}.

\begin{quote}
Das Kind hat Fieber, die Stirn ist ganz heiß.
\end{quote}

\begin{quote}
\emph{The child has a fever; its forehead is very hot.}
\end{quote}

Plural: \emph{die Stirnen}. \emph{Die Stirn runzeln} — to frown.

\begin{cousinweb}
\emph{Stirn} is inherited from Middle High German \emph{stirne} and Old High German \emph{stirna}, from Proto-Germanic \emph{stirnijō} \textquotedblleft{}forehead\textquotedblright{} (EtymDB).
\end{cousinweb}

\begin{grammarlens}[title={What you've built}]
One more word for telling, explaining and describing at B1.
\end{grammarlens}

\subsection*{Your turn}
\begin{itemize}
\raggedright
  \item \emph{Say it:} \emph{die Stirn}
  \item \emph{Say it:} \emph{die Stirn}, once more
  \item \emph{Say it:} a sentence of your own with \emph{die Stirn}
  \item \emph{Recall:} say the German for wheelchair, then the German for to heal, then say \emph{die Stirn} again
\end{itemize}

\begin{tcolorbox}[breakable,colback=teal!4,colframe=teal!35!black,title={Before you move on}]
What does die Stirn mean? (\textquotedblleft{}Forehead\textquotedblright{}.) Say it once more.
\end{tcolorbox}
\section[die Wange]{die Wange — cheek}

\label{lesson:GE-C484-wange}



\begin{quote}
Before the new one: say the German for to heal, then the German for forehead.
\end{quote}

\subsection*{You'll want to know: die Wange}
\textbf{die Wange} — \textquotedblleft{}cheek\textquotedblright{}.

\begin{quote}
Das Baby hat rote Wangen.
\end{quote}

\begin{quote}
\emph{The baby has red cheeks.}
\end{quote}

Plural: \emph{die Wangen}. In speech also \emph{die Backe}.

\begin{cousinweb}
\emph{Wange} is inherited from Middle High German \emph{wange} and Old High German \emph{wanga}, from Proto-Germanic \emph{wangô} \textquotedblleft{}cheek\textquotedblright{} (EtymDB).
\end{cousinweb}

\begin{grammarlens}[title={What you've built}]
One more word for telling, explaining and describing at B1.
\end{grammarlens}

\subsection*{Your turn}
\begin{itemize}
\raggedright
  \item \emph{Say it:} \emph{die Wange}
  \item \emph{Say it:} \emph{die Wange}, once more
  \item \emph{Say it:} a sentence of your own with \emph{die Wange}
  \item \emph{Recall:} say the German for to heal, then the German for forehead, then say \emph{die Wange} again
\end{itemize}

\begin{tcolorbox}[breakable,colback=teal!4,colframe=teal!35!black,title={Before you move on}]
What does die Wange mean? (\textquotedblleft{}Cheek\textquotedblright{}.) Say it once more.
\end{tcolorbox}
\section[das Kinn]{das Kinn — chin}

\label{lesson:GE-C484-kinn}



\begin{quote}
Before the new one: say the German for forehead, then the German for cheek.
\end{quote}

\subsection*{You'll want to know: das Kinn}
\textbf{das Kinn} — \textquotedblleft{}chin\textquotedblright{}.

\begin{quote}
Der Mann hat einen Bart am Kinn.
\end{quote}

\begin{quote}
\emph{The man has a beard on his chin.}
\end{quote}

Plural: \emph{die Kinne}.

\begin{cousinweb}
\emph{Kinn} is inherited from Old High German \emph{kinni}, from Proto-Germanic \emph{kinnuz} \textquotedblleft{}chin\textquotedblright{} (EtymDB); English \emph{chin} is a cousin.
\end{cousinweb}

\begin{grammarlens}[title={What you've built}]
One more word for telling, explaining and describing at B1.
\end{grammarlens}

\subsection*{Your turn}
\begin{itemize}
\raggedright
  \item \emph{Say it:} \emph{das Kinn}
  \item \emph{Say it:} \emph{das Kinn}, once more
  \item \emph{Say it:} a sentence of your own with \emph{das Kinn}
  \item \emph{Recall:} say the German for forehead, then the German for cheek, then say \emph{das Kinn} again
\end{itemize}

\begin{tcolorbox}[breakable,colback=teal!4,colframe=teal!35!black,title={Before you move on}]
What does das Kinn mean? (\textquotedblleft{}Chin\textquotedblright{}.) Say it once more.
\end{tcolorbox}
\section[die Augenbraue]{die Augenbraue — eyebrow}

\label{lesson:GE-C484-augenbraue}



\begin{quote}
Before the new one: say the German for cheek, then the German for chin.
\end{quote}

\subsection*{You'll want to know: die Augenbraue}
\textbf{die Augenbraue} — \textquotedblleft{}eyebrow\textquotedblright{}.

\begin{quote}
Überrascht hob sie die Augenbrauen.
\end{quote}

\begin{quote}
\emph{In surprise she raised her eyebrows.}
\end{quote}

Plural: \emph{die Augenbrauen}.

\begin{cousinweb}
\emph{Augenbraue} is inherited from Middle High German \emph{ougebrā}, a compound of \emph{ouge} \textquotedblleft{}eye\textquotedblright{} and \emph{brā} \textquotedblleft{}brow\textquotedblright{} (EtymDB).
\end{cousinweb}

\begin{grammarlens}[title={What you've built}]
One more word for telling, explaining and describing at B1.
\end{grammarlens}

\subsection*{Your turn}
\begin{itemize}
\raggedright
  \item \emph{Say it:} \emph{die Augenbraue}
  \item \emph{Say it:} \emph{die Augenbraue}, once more
  \item \emph{Say it:} a sentence of your own with \emph{die Augenbraue}
  \item \emph{Recall:} say the German for cheek, then the German for chin, then say \emph{die Augenbraue} again
\end{itemize}

\begin{tcolorbox}[breakable,colback=teal!4,colframe=teal!35!black,title={Before you move on}]
What does die Augenbraue mean? (\textquotedblleft{}Eyebrow\textquotedblright{}.) Say it once more.
\end{tcolorbox}
\section[die Wimper]{die Wimper — eyelash}

\label{lesson:GE-C484-wimper}



\begin{quote}
Before the new one: say the German for chin, then the German for eyebrow.
\end{quote}

\subsection*{You'll want to know: die Wimper}
\textbf{die Wimper} — \textquotedblleft{}eyelash\textquotedblright{}.

\begin{quote}
Das Mädchen hat lange, dunkle Wimpern.
\end{quote}

\begin{quote}
\emph{The girl has long, dark eyelashes.}
\end{quote}

Plural: \emph{die Wimpern}. \emph{Ohne mit der Wimper zu zucken} — without batting an eyelid.

\begin{cousinweb}
\emph{Wimper} is usually traced to Old High German \emph{wintbrāwa}, a compound with \emph{brāwa} \textquotedblleft{}brow\textquotedblright{}; the first part is not settled, and EtymDB lists the word without a source.
\end{cousinweb}

\begin{grammarlens}[title={What you've built}]
One more word for telling, explaining and describing at B1.
\end{grammarlens}

\subsection*{Your turn}
\begin{itemize}
\raggedright
  \item \emph{Say it:} \emph{die Wimper}
  \item \emph{Say it:} \emph{die Wimper}, once more
  \item \emph{Say it:} a sentence of your own with \emph{die Wimper}
  \item \emph{Recall:} say the German for chin, then the German for eyebrow, then say \emph{die Wimper} again
\end{itemize}

\begin{tcolorbox}[breakable,colback=teal!4,colframe=teal!35!black,title={Before you move on}]
What does die Wimper mean? (\textquotedblleft{}Eyelash\textquotedblright{}.) Say it once more.
\end{tcolorbox}
